%% toc: P15 159 159.8
\section{Sigma-algebra conditional independence}
\label{sec:annex-p15-w02c-p15-009-sigma-ci}

Throughout, \((\Omega,\mathcal F_0,P)\) is a probability space and all
sigma-algebras are sub-sigma-algebras of \(\mathcal F_0\). For an event
\(A\), a version of \(P(A\mid\mathcal H)\) is chosen with values in
\([0,1]\); this is possible by \cref{thm:p15-conditional-basic}, since
the version may be changed on the \(\mathcal H\)-measurable null set where
it leaves \([0,1]\). Equalities between conditional expectations hold
almost surely.

Two facts are used repeatedly.

\emph{Fact~F1 (finite measures agreeing on a generator).} Let \(\mu,\nu\)
be finite measures on a sigma-algebra \(\mathcal G=\sigma(\mathcal P)\),
where \(\mathcal P\) is a pi-system containing \(\Omega\), and let
\(\mu(E)=\nu(E)\) for all \(E\in\mathcal P\). Then \(\mu=\nu\) on
\(\mathcal G\). Indeed \(\mu(\Omega)=\nu(\Omega)=c\). If \(c=0\) both are
zero. Otherwise \(\mu/c\) and \(\nu/c\) are probability measures agreeing on
\(\mathcal P\), hence on \(\mathcal G\) by \cref{thm:p15-measure-uniqueness}.

\emph{Fact~F2 (the join is generated by intersections).} For
sigma-algebras \(\mathcal B,\mathcal H\), the collection
\(\mathcal P=\{B\cap H:B\in\mathcal B,H\in\mathcal H\}\) is a pi-system
containing \(\Omega\) and generating \(\mathcal B\vee\mathcal H\).
Indeed \((B\cap H)\cap(B'\cap H')=(B\cap B')\cap(H\cap H')\);
\(\Omega=\Omega\cap\Omega\); every member lies in \(\mathcal B\vee\mathcal H\);
and \(\mathcal P\) contains \(\mathcal B\) (take \(H=\Omega\)) and
\(\mathcal H\) (take \(B=\Omega\)), so \(\sigma(\mathcal P)\supseteq\mathcal B\vee\mathcal H\).

\begin{annexitem}{annex:p15-w02c-p15-009-sigma-ci-criteria}{thm:p15-w02c-p15-009-sigma-ci-criteria}{Criteria for conditional independence}
\begin{proof}
\paragraph{Item 1.} Fix \(C\in\mathcal C\) and let
\(\mathcal L_C=\{E\in\mathcal F_0:C\perp E\mid\mathcal H\}\). It is a
lambda-system (\cref{def:p15-pi-lambda}).
First, \(\Omega\in\mathcal L_C\): the constant \(1\) is
\(\mathcal H\)-measurable, so \(P(\Omega\mid\mathcal H)=1\) and
\(P(C\cap\Omega\mid\mathcal H)=P(C\mid\mathcal H)\cdot1\)
(\cref{thm:p15-conditional-basic}).
Second, if \(E\subseteq E'\) lie in \(\mathcal L_C\), then
\(\mathbf1_{C\cap(E'\setminus E)}=\mathbf1_{C\cap E'}-\mathbf1_{C\cap E}\) and
\(\mathbf1_{E'\setminus E}=\mathbf1_{E'}-\mathbf1_E\), so by linearity
\[
\begin{aligned}
 P(C\cap(E'\setminus E)\mid\mathcal H)
 &=P(C\mid\mathcal H)P(E'\mid\mathcal H)-P(C\mid\mathcal H)P(E\mid\mathcal H)\\
 &=P(C\mid\mathcal H)P(E'\setminus E\mid\mathcal H).
\end{aligned}
\]
Third, if \(E_n\in\mathcal L_C\) increase to \(E\), then
\(\mathbf1_{C\cap E_n}\uparrow\mathbf1_{C\cap E}\) and
\(\mathbf1_{E_n}\uparrow\mathbf1_E\), so by conditional monotone
convergence (\cref{thm:p15-conditional-limits}, item~1)
\(P(C\cap E_n\mid\mathcal H)\uparrow P(C\cap E\mid\mathcal H)\) and
\(P(E_n\mid\mathcal H)\uparrow P(E\mid\mathcal H)\) outside one null set.
Multiplying the second limit by the finite value \(P(C\mid\mathcal H)\)
and using the hypothesis for each \(n\) gives
\(P(C\cap E\mid\mathcal H)=P(C\mid\mathcal H)P(E\mid\mathcal H)\).
Since \(\mathcal D\subseteq\mathcal L_C\) and \(\mathcal D\) is a
pi-system, \cref{thm:p15-measure-uniqueness} gives
\(\sigma(\mathcal D)\subseteq\mathcal L_C\). Thus
\(\mathcal C\perp\sigma(\mathcal D)\mid\mathcal H\).

Now fix \(D\in\sigma(\mathcal D)\) and let
\(\mathcal L'_D=\{E\in\mathcal F_0:E\perp D\mid\mathcal H\}\). The relation
for two events is symmetric, since \(E\cap D=D\cap E\) and real
multiplication commutes, so \(\mathcal L'_D=\mathcal L_D\) is a
lambda-system by the argument just given. By the previous paragraph it
contains \(\mathcal C\), hence \(\sigma(\mathcal C)\). Thus
\(\sigma(\mathcal C)\perp\sigma(\mathcal D)\mid\mathcal H\).

\paragraph{Item 2.} \emph{Simple variables.} If
\(X=\sum_i a_i\mathbf1_{A_i}\) and \(Y=\sum_j b_j\mathbf1_{B_j}\) are finite
sums with \(A_i\in\mathcal A\), \(B_j\in\mathcal B\), then
\(XY=\sum_{i,j}a_ib_j\mathbf1_{A_i\cap B_j}\), and linearity with the
definition gives
\[
 E(XY\mid\mathcal H)=\sum_{i,j}a_ib_jP(A_i\mid\mathcal H)P(B_j\mid\mathcal H)
 =E(X\mid\mathcal H)E(Y\mid\mathcal H).
\]
\emph{Nonnegative variables.} Let \(X,Y\geq0\) satisfy the hypotheses. Put
\(X_n=\min(n,2^{-n}\lfloor2^nX\rfloor)\) and define \(Y_n\) likewise. Each
\(X_n\) is a simple \(\mathcal A\)-measurable variable, each \(Y_n\) a
simple \(\mathcal B\)-measurable one, \(0\leq X_n\uparrow X\),
\(0\leq Y_n\uparrow Y\), and hence \(0\leq X_nY_n\uparrow XY\). By the
simple case \(E(X_nY_n\mid\mathcal H)=E(X_n\mid\mathcal H)E(Y_n\mid\mathcal H)\)
for every \(n\). By \cref{thm:p15-conditional-limits}, item~1, outside one
null set the three sequences increase to \(E(XY\mid\mathcal H)\),
\(E(X\mid\mathcal H)\) and \(E(Y\mid\mathcal H)\), which are finite almost
surely because \(XY,X,Y\) are integrable. Passing to the limit in the
product of the two finite nondecreasing sequences gives the identity.

\emph{Signed variables.} Write \(X=X^+-X^-\) and \(Y=Y^+-Y^-\). The parts
\(X^\pm\) are \(\mathcal A\)-measurable and integrable, \(Y^\pm\) are
\(\mathcal B\)-measurable and integrable, and each product
\(X^aY^b\) (\(a,b\in\{+,-\}\)) satisfies \(|X^aY^b|\leq|XY|\), so it is
integrable. The nonnegative case applies to each pair. Since
\(XY=X^+Y^+-X^+Y^--X^-Y^++X^-Y^-\), linearity gives
\[
\begin{aligned}
 E(XY\mid\mathcal H)
 &=\bigl(E(X^+\mid\mathcal H)-E(X^-\mid\mathcal H)\bigr)
  \bigl(E(Y^+\mid\mathcal H)-E(Y^-\mid\mathcal H)\bigr)\\
 &=E(X\mid\mathcal H)E(Y\mid\mathcal H).
\end{aligned}
\]

\paragraph{A common identity.} For \(A\in\mathcal A\), \(B\in\mathcal B\)
and \(H\in\mathcal H\),
\begin{equation}
\begin{aligned}
 \int_{B\cap H}P(A\mid\mathcal H)\,dP
 &=\int_H\mathbf1_BP(A\mid\mathcal H)\,dP
 =\int_H E\bigl(\mathbf1_BP(A\mid\mathcal H)\bigm|\mathcal H\bigr)\,dP\\
 &=\int_H P(A\mid\mathcal H)P(B\mid\mathcal H)\,dP .
\end{aligned}
 \tag{\(*\)}
\end{equation}
The second equality is the defining identity of conditional expectation on
\(H\in\mathcal H\) (\cref{def:p15-general-conditional-expectation}) for the
bounded variable \(\mathbf1_BP(A\mid\mathcal H)\). The third takes out the
bounded \(\mathcal H\)-measurable factor \(P(A\mid\mathcal H)\)
(\cref{thm:p15-conditional-pullout}).

\paragraph{Item 3.} Suppose \(\mathcal A\perp\mathcal B\mid\mathcal H\) and
fix \(A\in\mathcal A\). By the hypothesis, the right side of \((*)\) equals
\(\int_HP(A\cap B\mid\mathcal H)\,dP=P(A\cap B\cap H)\). Hence the finite
measures \(E\mapsto\int_EP(A\mid\mathcal H)\,dP\) and
\(E\mapsto P(A\cap E)\) on \(\mathcal B\vee\mathcal H\) agree on the
pi-system of Fact~F2, and so agree on \(\mathcal B\vee\mathcal H\) by
Fact~F1. As \(P(A\mid\mathcal H)\) is \(\mathcal H\)-measurable, hence
\(\mathcal B\vee\mathcal H\)-measurable, it is a version of
\(P(A\mid\mathcal B\vee\mathcal H)\), and by uniqueness of versions
(\cref{thm:p15-conditional-existence}) every version equals it almost
surely.

Conversely, suppose \(P(A\mid\mathcal B\vee\mathcal H)=P(A\mid\mathcal H)\)
for every \(A\in\mathcal A\). Fix \(A\in\mathcal A\) and \(B\in\mathcal B\).
For \(H\in\mathcal H\), the event \(B\cap H\) lies in \(\mathcal B\vee\mathcal H\),
so the left side of \((*)\) equals
\(\int_{B\cap H}P(A\mid\mathcal B\vee\mathcal H)\,dP=P(A\cap B\cap H)\).
Hence \(\int_HP(A\mid\mathcal H)P(B\mid\mathcal H)\,dP=P(A\cap B\cap H)\)
for all \(H\in\mathcal H\). The product is bounded and
\(\mathcal H\)-measurable, so it is a version of \(P(A\cap B\mid\mathcal H)\).
Thus \(A\perp B\mid\mathcal H\) for all such \(A,B\).

\paragraph{Item 4.} If \(\mathcal A\perp\mathcal B\mid\mathcal H\), item~3
shows that \(Z_A=P(A\mid\mathcal H)\) serves. Conversely, let \(Z_A\) be
\(\mathcal H\)-measurable and a version of \(P(A\mid\mathcal B\vee\mathcal H)\).
For \(H\in\mathcal H\subseteq\mathcal B\vee\mathcal H\),
\(\int_HZ_A\,dP=P(A\cap H)\), so \(Z_A\) is also a version of
\(P(A\mid\mathcal H)\). Hence
\(P(A\mid\mathcal B\vee\mathcal H)=Z_A=P(A\mid\mathcal H)\) for every
\(A\in\mathcal A\), and item~3 gives \(\mathcal A\perp\mathcal B\mid\mathcal H\).

\paragraph{Item 5.} First let \(X\geq0\) be integrable and
\(\mathcal A\)-measurable. For \(B\in\mathcal B\) and \(H\in\mathcal H\),
\[
\begin{aligned}
 \int_{B\cap H}X\,dP&=\int_HE(\mathbf1_BX\mid\mathcal H)\,dP
 =\int_HP(B\mid\mathcal H)E(X\mid\mathcal H)\,dP\\
 &=\int_HE\bigl(\mathbf1_BE(X\mid\mathcal H)\bigm|\mathcal H\bigr)\,dP
 =\int_{B\cap H}E(X\mid\mathcal H)\,dP .
\end{aligned}
\]
The first and last equalities are the defining identity on \(H\). The
second is item~2 with the \(\mathcal B\)-measurable factor \(\mathbf1_B\),
for which \(E|\mathbf1_BX|\leq E|X|<\infty\). The third takes out the
\(\mathcal H\)-measurable factor \(E(X\mid\mathcal H)\), which is
integrable, from the integrable product \(\mathbf1_BE(X\mid\mathcal H)\)
(\cref{thm:p15-conditional-pullout}). Both sides are finite measures in the
event (the version \(E(X\mid\mathcal H)\) may be taken nonnegative by
\cref{thm:p15-conditional-basic}), and they agree on the pi-system of
Fact~F2, hence on \(\mathcal B\vee\mathcal H\) by Fact~F1. So the
\(\mathcal H\)-measurable \(E(X\mid\mathcal H)\) is a version of
\(E(X\mid\mathcal B\vee\mathcal H)\). For a signed integrable \(X\), apply
this to \(X^+\) and \(X^-\) and use linearity.

\paragraph{Item 6.} If \(\mathcal B\subseteq\mathcal H\), then
\(\mathcal B\vee\mathcal H=\mathcal H\), so the criterion of item~3 holds
with equality of the same object.

Let every event of \(\mathcal H\) have probability \(0\) or \(1\). For any
event \(A\) the constant \(P(A)\) is \(\mathcal H\)-measurable and, for
\(H\in\mathcal H\), \(\int_HP(A)\,dP=P(A)P(H)=P(A\cap H)\): if \(P(H)=0\)
both sides vanish, and if \(P(H)=1\) then
\(P(A\cap H)=P(A)-P(A\cap H^c)=P(A)\). So \(P(A\mid\mathcal H)=P(A)\) for
every event. The defining identity of the relation becomes
\(P(A\cap B)=P(A)P(B)\) almost surely, and an almost-sure equality of two
constants is an equality of numbers.

If \(\mathcal H=\mathcal F_0\), every event \(F\) has the
\(\mathcal F_0\)-measurable version \(P(F\mid\mathcal F_0)=\mathbf1_F\)
(\cref{thm:p15-conditional-basic}), and
\(\mathbf1_A\mathbf1_B=\mathbf1_{A\cap B}\).

\paragraph{Item 7.} Every event of \(\sigma(X)\) has the form
\(\{X\in D\}\) for a measurable set \(D\) of the value space, and similarly
for \(Y\). If \(X\perp Y\mid Z\) and \(u,v\) are bounded measurable, then
\(u(X)\) is \(\sigma(X)\)-measurable, \(v(Y)\) is \(\sigma(Y)\)-measurable,
and both and their product are bounded, so item~2 with
\(\mathcal H=\sigma(Z)\) gives the product identity. Conversely, taking
\(u=\mathbf1_D\) and \(v=\mathbf1_F\) in the product identity gives
\(\{X\in D\}\perp\{Y\in F\}\mid\sigma(Z)\) for all measurable \(D,F\), which
is \(\sigma(X)\perp\sigma(Y)\mid\sigma(Z)\). For standard Borel spaces the
product identity is the defining condition of \cref{def:p16-measurable-ci},
in which \(E[\cdot\mid Z]\) also means conditioning on \(\sigma(Z)\).
\end{proof}
\end{annexitem}

\begin{annexitem}{annex:p15-w02c-p15-009-sigma-ci-rules}{thm:p15-w02c-p15-009-sigma-ci-rules}{The inference rules}
\begin{proof}
\paragraph{Symmetry and reduction.} The defining identity for events
\(A,B\) is unchanged when they are exchanged, because
\(A\cap B=B\cap A\) and real multiplication commutes; so the relation for
classes is symmetric. If \(\mathcal D\subseteq\mathcal B\), every pair
\((A,D)\) with \(A\in\mathcal A\), \(D\in\mathcal D\) is a pair of
\(\mathcal A\times\mathcal B\), so reduction holds. Decomposition is
reduction applied to \(\mathcal B\subseteq\mathcal B\vee\mathcal D\) and to
\(\mathcal D\subseteq\mathcal B\vee\mathcal D\).

\paragraph{Weak union.} Let \(\mathcal A\perp(\mathcal B\vee\mathcal C)\mid\mathcal D\)
and \(A\in\mathcal A\). Item~3 of
\cref{thm:p15-w02c-p15-009-sigma-ci-criteria}, with the join taken in
either order since \(\vee\) is associative and commutative, gives
\(P(A\mid\mathcal B\vee\mathcal C\vee\mathcal D)=P(A\mid\mathcal D)\).
The version \(P(A\mid\mathcal D)\) is \(\mathcal D\)-measurable, hence
\(\mathcal C\vee\mathcal D\)-measurable, and it is a version of the
conditional probability given the larger sigma-algebra
\(\mathcal B\vee\mathcal C\vee\mathcal D\supseteq\mathcal C\vee\mathcal D\).
By the last sentence of \cref{thm:p15-conditional-tower}, it therefore
represents \(P(A\mid\mathcal C\vee\mathcal D)\). So
\(P(A\mid\mathcal B\vee(\mathcal C\vee\mathcal D))=P(A\mid\mathcal C\vee\mathcal D)\)
for every \(A\in\mathcal A\), and item~3 with
\(\mathcal H=\mathcal C\vee\mathcal D\) gives
\(\mathcal A\perp\mathcal B\mid\mathcal C\vee\mathcal D\).

\paragraph{Contraction.} Let \(\mathcal A\perp\mathcal C\mid\mathcal B\),
\(\mathcal A\perp\mathcal D\mid\mathcal B\vee\mathcal C\) and
\(A\in\mathcal A\). Item~3 applied to the second premise, with
\(\mathcal H=\mathcal B\vee\mathcal C\), and then to the first, with
\(\mathcal H=\mathcal B\), gives
\[
 P(A\mid\mathcal D\vee\mathcal B\vee\mathcal C)=P(A\mid\mathcal B\vee\mathcal C)
 =P(A\mid\mathcal B).
\]
Since \(\mathcal D\vee\mathcal B\vee\mathcal C=(\mathcal C\vee\mathcal D)\vee\mathcal B\),
item~3 in the converse direction gives
\(\mathcal A\perp(\mathcal C\vee\mathcal D)\mid\mathcal B\).

\paragraph{Random elements.} Let \(Y,Z\) take values in
\((S_Y,\mathcal S_Y)\) and \((S_Z,\mathcal S_Z)\), and give
\(S_Y\times S_Z\) the product sigma-algebra, generated by the rectangles
\(D\times F\). The class of sets \(G\) of the product sigma-algebra with
\(\{(Y,Z)\in G\}\in\sigma(Y)\vee\sigma(Z)\) is a sigma-algebra (preimages
commute with complements and countable unions) containing every rectangle,
since \(\{(Y,Z)\in D\times F\}=\{Y\in D\}\cap\{Z\in F\}\); hence it is the
whole product sigma-algebra and \(\sigma((Y,Z))\subseteq\sigma(Y)\vee\sigma(Z)\).
Conversely \(\{Y\in D\}=\{(Y,Z)\in D\times S_Z\}\) and
\(\{Z\in F\}=\{(Y,Z)\in S_Y\times F\}\) lie in \(\sigma((Y,Z))\), so
\(\sigma(Y)\vee\sigma(Z)\subseteq\sigma((Y,Z))\). For a measurable
\(\varphi\), \(\{\varphi(Y)\in E\}=\{Y\in\varphi^{-1}(E)\}\in\sigma(Y)\).
With these identities, the four statements for random elements are items~1,
2, 3 and~4 applied to \(\sigma(X),\sigma(Y),\sigma(Z),\sigma(W)\).
\end{proof}
\end{annexitem}

\begin{annexitem}{annex:p15-w02c-p15-009-ci-tables}{thm:p15-w02c-p15-009-ci-tables}{Created and destroyed independence; intersection}
\begin{proof}
For \(\mathcal H=\{\varnothing,C,C^c,\Omega\}\) with \(0<P(C)<1\),
\cref{thm:p15-conditional-partitions} gives the version of
\(P(E\mid\mathcal H)\) equal to \(P(E\cap C)/P(C)\) on \(C\) and to
\(P(E\cap C^c)/P(C^c)\) on \(C^c\). Two such versions agree almost surely
exactly when the two constants agree on each cell, since both cells have
positive probability. So \(A\perp B\mid\mathcal H\) holds exactly when
\(P(A\cap B\mid C)=P(A\mid C)P(B\mid C)\) and the same holds with \(C^c\),
where \(P(\cdot\mid C)=P(\cdot\cap C)/P(C)\).

\paragraph{Item 1.} The cells on \(C\) sum to \(9/18\) and those on \(C^c\)
also to \(9/18\). On \(C\): \(P(A\mid C)=3/9\), \(P(B\mid C)=6/9\) and
\(P(A\cap B\mid C)=2/9=(3/9)(6/9)\). On \(C^c\): \(P(A\mid C^c)=6/9\),
\(P(B\mid C^c)=3/9\) and \(P(A\cap B\mid C^c)=2/9=(6/9)(3/9)\). Hence
\(A\perp B\mid\mathcal H\). Unconditionally \(P(A)=(2+1+2+4)/18=1/2\),
\(P(B)=(2+4+2+1)/18=1/2\) and \(P(A\cap B)=(2+2)/18=2/9\ne1/4\).

\paragraph{Item 2.} Here \(P(A)=(1+2+2+1)/12=1/2\),
\(P(B)=(1+2+2+1)/12=1/2\) and \(P(A\cap B)=(1+2)/12=1/4=P(A)P(B)\), so
\(A\) and \(B\) are independent. The cells on \(C\) sum to \(6/12\), and
\(P(A\mid C)=3/6\), \(P(B\mid C)=3/6\), \(P(A\cap B\mid C)=1/6\ne1/4\).
So the condition fails on \(C\), and \(A\perp B\mid\mathcal H\) fails.

\paragraph{Item 3.} Let \(N=\{(0,0,0),(1,1,1)\}^c\), a null set. Off \(N\)
the three coordinates are equal, so \(X_1=X_2=X_3\) almost surely. Each
\(X_i\) takes the values \(0\) and \(1\) with probability \(1/2\), and
\(\mathcal A_i=\{\varnothing,\{X_i=0\},\{X_i=1\},\Omega\}\).

\emph{The two premises.} Let \(A\in\mathcal A_1\); then
\(\mathbf1_A=g(X_1)\) for some \(g:\{0,1\}\to\{0,1\}\). The variable
\(g(X_3)\) is \(\mathcal A_3\)-measurable, hence
\(\mathcal A_2\vee\mathcal A_3\)-measurable, and equals \(\mathbf1_A\)
off \(N\). So \(\int_Gg(X_3)\,dP=P(A\cap G)\) for every
\(G\in\mathcal A_2\vee\mathcal A_3\), and \(g(X_3)\) is an
\(\mathcal A_3\)-measurable version of \(P(A\mid\mathcal A_2\vee\mathcal A_3)\).
Item~4 of \cref{thm:p15-w02c-p15-009-sigma-ci-criteria} gives
\(\mathcal A_1\perp\mathcal A_2\mid\mathcal A_3\). Exchanging the roles of
\(X_2\) and \(X_3\) gives \(\mathcal A_1\perp\mathcal A_3\mid\mathcal A_2\).

\emph{The intersection.} The point \((0,1,0)\) lies in \(\{X_2=1\}\) but not
in \(\{X_3=1\}\), and \((0,0,1)\) lies in \(\{X_2=0\}\) but not in
\(\{X_3=0\}\); and \(\{X_2=0\}\ne\{X_3=1\}\) because \((0,0,0)\) lies in the
first only. So no nonempty proper member of \(\mathcal A_2\) belongs to
\(\mathcal A_3\), and \(\mathcal A_2\cap\mathcal A_3=\{\varnothing,\Omega\}\).

\emph{The failed conclusion.} Given the trivial sigma-algebra, item~6 of
\cref{thm:p15-w02c-p15-009-sigma-ci-criteria} reduces
\(\mathcal A_1\perp(\mathcal A_2\vee\mathcal A_3)\mid\{\varnothing,\Omega\}\)
to independence. But \(\{X_1=1\}\in\mathcal A_1\) and
\(\{X_2=1\}\in\mathcal A_2\vee\mathcal A_3\) satisfy
\(P(X_1=1,X_2=1)=P\{(1,1,1)\}=1/2\ne1/4=P(X_1=1)P(X_2=1)\).
\end{proof}
\end{annexitem}
